\documentclass[11pt,a4paper]{article}
\usepackage[utf8]{inputenc}
\usepackage[T1]{fontenc}
\usepackage[spanish,es-nodecimaldot]{babel}
\usepackage{lmodern}
\usepackage{amsmath}
\usepackage{booktabs}
\usepackage{geometry}
\usepackage{xcolor}
\usepackage{hyperref}
\usepackage{listings}
\usepackage{microtype}
\geometry{margin=2.35cm}
\hypersetup{colorlinks=true,linkcolor=blue!55!black,urlcolor=blue!55!black}
\lstdefinestyle{terminal}{basicstyle=\ttfamily\small,breaklines=true,
  frame=single,showstringspaces=false,columns=fullflexible}

\title{¿Qué está haciendo ROS 2 actualmente en el proyecto?\\Nova Spot Micro}
\author{Proyecto de tesis --- robot cuadrúpedo Nova Spot Micro}
\date{Actualizado el 16 de agosto de 2026}

\begin{document}
\maketitle
\begin{abstract}
Este documento explica el papel que cumple ROS 2 Humble en el estado actual del
proyecto Nova Spot Micro. Describe cómo se carga el modelo, cómo se inicia
Gazebo o MuJoCo, cómo se generan y transmiten las referencias articulares, cómo
se observan los estados y cuáles componentes faltan para llegar al robot físico
con Raspberry Pi 3, PCA9685, MG996R y sensores. El sistema actual constituye
una arquitectura funcional de simulación, no una interfaz de hardware terminada.
\end{abstract}
\tableofcontents

\section{Papel de ROS 2 en el proyecto}
ROS 2 es la infraestructura que conecta las partes del sistema. No calcula por
sí solo cómo caminar ni simula la física: coordina nodos, parámetros, mensajes,
tiempo y controladores para que el código de marcha, el modelo y el simulador
trabajen como un solo sistema.

En el estado actual, ROS 2 realiza además observación y seguridad:
\begin{enumerate}
\item procesa y distribuye la descripción del robot;
\item inicia los procesos necesarios mediante archivos \texttt{launch};
\item recibe comandos de alto nivel como \texttt{gateo} o \texttt{stand};
\item transporta referencias para las doce articulaciones mediante
  \texttt{JointTrajectory};
\item publica los estados articulares producidos por el simulador.
\item publica fase, métricas de pose y contactos de los cuatro pies;
\item compara contacto previsto y observado y supervisa altura e inclinación.
\end{enumerate}

\section{Paquetes ROS 2 desarrollados}
El workspace contiene dos paquetes propios importantes:
\begin{table}[h]
\centering
\begin{tabular}{p{0.34\textwidth}p{0.58\textwidth}}
\toprule
Paquete & Función\\
\midrule
\texttt{nova\_sm3\_description} & Contiene URDF/Xacro, modelo MJCF,
configuración de \texttt{ros2\_control}, RViz y lanzadores de Gazebo/MuJoCo.\\
\texttt{nova\_gait\_controller} & Contiene cinemática directa e inversa,
generación de gateo y marcha paso, máquina de estados, métricas, supervisor,
contactos, parámetros y lanzadores de demostración.\\
\bottomrule
\end{tabular}
\end{table}

Estos paquetes se compilan con \texttt{colcon}. La opción
\texttt{--symlink-install} permite que los cambios en archivos Python y datos
instalados se reflejen con mayor facilidad durante el desarrollo.

\section{Descripción del robot y \texttt{robot\_description}}
El archivo Xacro describe el cuerpo, cuatro patas, doce articulaciones, enlaces,
colisiones, masas e inercias provisionales. También declara nombres, ejes,
límites y las interfaces que utilizará \texttt{ros2\_control}.

Al iniciar la simulación, ROS 2 ejecuta \texttt{xacro} y convierte las macros en
URDF. El texto completo se entrega como el parámetro
\texttt{robot\_description}. A partir de ese parámetro:
\begin{itemize}
\item \texttt{robot\_state\_publisher} conoce la estructura cinemática;
\item Gazebo crea el modelo físico;
\item el complemento de \texttt{ros2\_control} descubre las doce articulaciones;
\item RViz puede representar enlaces y transformaciones.
\end{itemize}

Las posiciones iniciales articulares son
$(0.10,0.42,-0.84)$ rad por pata. Esto permite que el robot aparezca en postura
en lugar de hacerlo con las piernas rectas mientras arrancan los controladores.

\section{Procesos que se ejecutan en Gazebo}
El comando principal es
\begin{lstlisting}[style=terminal]
ros2 launch nova_gait_controller demo.launch.py
\end{lstlisting}
Este lanzador incluye \texttt{sim.launch.py} y, después de cinco segundos,
inicia el nodo \texttt{gait\_controller}. El conjunto resultante contiene:

\subsection{Gazebo Sim}
Calcula gravedad, contactos, colisiones y movimiento del cuerpo. El robot se
crea con nombre \texttt{nova\_sm3} a una altura inicial de $0.245$ m. Gazebo
no decide los pasos: recibe las órdenes articulares desde ROS 2 y devuelve el
estado físico resultante.

\subsection{\texttt{robot\_state\_publisher}}
Publica las transformaciones entre enlaces usando el URDF y los estados
articulares. Gracias a él, herramientas como RViz pueden reconstruir la postura
completa a partir de doce ángulos.

\subsection{\texttt{controller\_manager}}
Administra interfaces y ciclos de control. En Gazebo está incorporado mediante
\texttt{gz\_ros2\_control}; en MuJoCo se ejecuta como
\texttt{ros2\_control\_node} con la interfaz específica de MuJoCo.

\subsection{Spawners de controladores}
Dos procesos temporales solicitan al administrador cargar, configurar y activar:
\begin{itemize}
\item \texttt{joint\_state\_broadcaster};
\item \texttt{joint\_trajectory\_controller}.
\end{itemize}
Cuando terminan exitosamente, desaparecen; los controladores permanecen activos
dentro del administrador.

\subsection{Nodo \texttt{gait\_controller}}
Es el código propio que recibe el modo de marcha, genera la tabla cartesiana,
aplica cinemática inversa y publica un punto articular en cada fase.

\section{Comandos y temas ROS 2}
El usuario selecciona el comportamiento publicando un mensaje de texto:
\begin{lstlisting}[style=terminal]
ros2 topic pub --once /nova/gait_command \
  std_msgs/msg/String "{data: gateo}"
\end{lstlisting}
El nodo acepta \texttt{stand}, \texttt{gateo/crawl}, \texttt{paso/step},
\texttt{galope/gallop} y \texttt{stop/parar}. Para la tesis, el gateo es la
marcha nominal desarrollada; el galope queda como demostración opcional de
simulación, deshabilitada por defecto mediante
\texttt{enable\_experimental\_gallop=false} y excluida del hardware.

Los temas principales son:
\begin{table}[h]
\centering
\small
\begin{tabular}{p{0.39\textwidth}p{0.20\textwidth}p{0.33\textwidth}}
\toprule
Tema & Tipo & Uso\\
\midrule
\texttt{/nova/gait\_command} & \texttt{std\_msgs/String} & Comando de alto nivel.\\
\texttt{/joint\_trajectory\_controller/joint\_trajectory} &
\texttt{trajectory\_msgs/JointTrajectory} & Referencias para 12 articulaciones.\\
\texttt{/joint\_states} & \texttt{sensor\_msgs/JointState} & Posición y velocidad simuladas.\\
\texttt{/nova/gait\_phase} & \texttt{std\_msgs/String} & Fase y contactos previstos.\\
\texttt{/nova/metrics/json} & \texttt{std\_msgs/String} & Pose y articulaciones sincronizadas.\\
\texttt{/nova/foot\_contacts} & \texttt{std\_msgs/String} & Contactos medidos agregados.\\
\texttt{/nova/contact\_diagnostics} & \texttt{std\_msgs/String} & Comparación previsto--observado.\\
\texttt{/robot\_description} o parámetro equivalente & URDF/XML & Descripción completa.\\
\bottomrule
\end{tabular}
\end{table}

\section{Cómo fluye una orden de gateo}
Cuando se publica \texttt{gateo}, ocurre esta secuencia:
\begin{enumerate}
\item DDS/ROS 2 entrega el mensaje al suscriptor de
  \texttt{gait\_controller}.
\item El nodo lee los parámetros: 24 muestras, paso de 18 mm, elevación de
  14 mm y 0.18 s por muestra.
\item Se valida que los valores estén dentro de los intervalos permitidos.
\item Se generan posiciones cartesianas para los cuatro pies. Tres permanecen
  en apoyo y una entra en oscilación.
\item La cinemática inversa convierte cada posición de pie en coxa, fémur y
  tibia, y rechaza puntos inalcanzables o fuera de límites.
\item Cada vector de doce ángulos se envía como \texttt{JointTrajectory}.
\item \texttt{joint\_trajectory\_controller} interpola la referencia.
\item \texttt{gz\_ros2\_control} aplica los comandos a las articulaciones de
  Gazebo.
\item Gazebo resuelve física y contactos.
\item \texttt{joint\_state\_broadcaster} publica el estado resultante.
\end{enumerate}

El ciclo nominal dura
\[
24\text{ muestras}\times0.18\text{ s}=4.32\text{ s}.
\]

\section{Función de \texttt{ros2\_control}}
\texttt{ros2\_control} separa el algoritmo de marcha del destino físico o
simulado. El controlador publica el mismo concepto de posición articular,
mientras una interfaz de sistema decide cómo aplicarlo.

Actualmente existen estos destinos:
\begin{itemize}
\item \textbf{hardware falso}: útil para verificar mensajes sin dinámica;
\item \textbf{Gazebo}: \texttt{gz\_ros2\_control/GazeboSimSystem};
\item \textbf{MuJoCo}: \texttt{mujoco\_ros2\_control/MujocoSystemInterface}.
\end{itemize}

No existe todavía una interfaz física para PCA9685. Por tanto, seleccionar una
marcha hoy mueve la simulación, no los MG996R.

\section{RViz, Gazebo y MuJoCo}
Las tres herramientas cumplen tareas distintas:
\begin{itemize}
\item \textbf{RViz}: inspecciona geometría, ejes, transformaciones y postura;
  no simula gravedad ni contacto.
\item \textbf{Gazebo}: verifica integración URDF/ROS 2, controladores, colisión,
  gravedad y marcha nominal.
\item \textbf{MuJoCo}: será el entorno principal de entrenamiento; hoy valida
  cadencia y articulaciones, pero no expone pose y contactos equivalentes.
\end{itemize}

El lanzador \texttt{mujoco\_demo.launch.py} usa el mismo
\texttt{gait\_controller} y el mismo YAML; cambia la interfaz de simulación.
Esto ayuda a evitar que la lógica de marcha dependa de un único simulador.

\section{Parámetros ROS 2 actuales}
El archivo \texttt{gaits.yaml} configura el nodo:
\begin{lstlisting}[style=terminal]
initial_gait: stand
crawl_phase_duration: 0.18
crawl_samples: 24
crawl_step_length: 0.018
crawl_step_height: 0.014
transition_ratio: 0.90
command_topic: /nova/gait_command
phase_topic: /nova/gait_phase
trajectory_topic: /joint_trajectory_controller/joint_trajectory
\end{lstlisting}
Los parámetros se cargan al iniciar el nodo. El gateo se vuelve a construir al
recibir el comando, de modo que se validan los valores vigentes antes de mover.

\section{Qué se ha verificado con ROS 2}
Hasta el momento se comprobó:
\begin{itemize}
\item compilación de los dos paquetes con \texttt{colcon};
\item carga del modelo y de las doce interfaces articulares;
\item activación de ambos controladores;
\item recepción de comandos de marcha;
\item publicación y seguimiento de trayectorias de doce articulaciones;
\item movimiento observable en \texttt{/joint\_states};
\item gateo con cinemática inversa en Gazebo durante más de diez ciclos;
\item avance aproximado de $0.219$ m sin caída en esa prueba;
\item altura entre $0.22287$ y $0.22428$ m;
\item ejecución previa del gateo en MuJoCo con estados articulares.
\item marcha paso reproducida en Gazebo y validada articularmente en MuJoCo;
\item cadencia corregida de 4.32 s por ciclo de gateo;
\item métricas sincronizadas, supervisor y cuatro contactos en Gazebo;
\item ensayo de 76 ciclos que detectó deslizamiento trasero.
\end{itemize}

ROS 2 permitió mantener la marcha, el simulador y el observador desacoplados,
así como repetir las pruebas mediante comandos explícitos.

\section{Raspberry Pi 3 y comunicación futura}
Se preparó el script \texttt{preparar\_raspberry\_pi\_ros2.sh} para Ubuntu
Server 22.04 ARM64. Instalará ROS 2 Humble base, SSH, I2C,
\texttt{ros2\_control} y controladores. También configura
\texttt{ROS\_DOMAIN\_ID=30} y \texttt{ROS\_LOCALHOST\_ONLY=0}.

La arquitectura prevista es:
\begin{itemize}
\item el computador ejecuta desarrollo, Gazebo, MuJoCo y entrenamiento;
\item la Raspberry Pi ejecuta control en tiempo real práctico, sensores,
  interfaz PCA9685 y registro básico;
\item ambos equipos comparten red y el mismo \texttt{ROS\_DOMAIN\_ID};
\item SSH permite administrar la Pi, mientras DDS transporta temas ROS 2.
\end{itemize}

Raspberry Pi Imager 2.0.10 ya fue descargado, pero no se ha grabado la microSD
porque todavía no se conectó una. Tampoco se ha detectado una Pi en la red.

\section{Qué todavía no hace ROS 2 en este proyecto}
Es importante no confundir simulación funcional con robot físico terminado.
Actualmente ROS 2 no:
\begin{itemize}
\item genera PWM para PCA9685;
\item conoce centros, sentidos y límites físicos de los doce MG996R;
\item mide corriente, tensión o temperatura;
\item recibe todavía BNO055 ni contactos de pies reales;
\item conoce el estado de una parada física de emergencia;
\item calcula margen de estabilidad en línea;
\item detecta contacto corporal o caída mediante una señal dedicada; el
  supervisor actual actúa por altura e inclinación;
\item ejecuta una política de aprendizaje por refuerzo;
\item garantiza comportamiento del robot real a partir de la simulación.
\end{itemize}

El comando lógico \texttt{stop} conserva postura en el controlador; no corta
energía. La parada física debe actuar independientemente de ROS 2.

\section{Próxima arquitectura ROS 2}
Los siguientes componentes pendientes son:
\begin{enumerate}
\item corregir la trayectoria de gateo y repetir la medición de contactos;
\item incorporar IMU y contactos equivalentes en MuJoCo;
\item calcular polígono de soporte y margen de estabilidad en línea;
\item ampliar el supervisor a límites, pérdida de datos y referencias inválidas,
  con mecanismo explícito de rearme;
\item una interfaz \texttt{ros2\_control} para PCA9685 con conversiones
  radianes--PWM calibradas;
\item nodos de BNO055, corriente/tensión, contactos y parada física;
\item registro mediante \texttt{rosbag2} para comparar marcha nominal y marcha
  con corrección aprendida.
\end{enumerate}

La política de aprendizaje por refuerzo se añadirá después de validar estos
elementos. Su salida será una corrección pequeña y acotada; ROS 2 seguirá
transportando referencias, pero el supervisor conservará la autoridad final.

\section{Comandos útiles de inspección}
\begin{lstlisting}[style=terminal]
# Ver nodos
ros2 node list

# Ver temas y tipos
ros2 topic list -t

# Ver controladores activos
ros2 control list_controllers

# Observar una muestra articular
ros2 topic echo /joint_states --once

# Ver parametros del nodo de marcha
ros2 param list /gait_controller

# Enviar postura segura logica
ros2 topic pub --once /nova/gait_command \
  std_msgs/msg/String "{data: stand}"
\end{lstlisting}

\section{Resumen}
ROS 2 ya integra el modelo, los dos simuladores, dos marchas nominales, la
cinemática inversa, métricas, fase, contactos de Gazebo y un supervisor
provisional. El siguiente salto es corregir el patrón físico de apoyo, añadir
estabilidad en línea y equivalencia sensorial en MuJoCo. La transferencia a
servos sigue bloqueada hasta caracterizar, proteger y calibrar el hardware.

\end{document}
